\chapter{Deducción de la ecuación implícita del dioptrio estigmático}
\label{apx:ovoid_implicit_deduction}

Este apéndice deduce, a partir de la condición estigmática de Fermat, la
ecuación polinómica implícita que caracteriza a la superficie cartesiana. La
forma canónica y su expansión en series se recogen en el
apéndice~\ref{apx:ovoids_series}.

\medskip

Sean $A=(0,0,z_0)$ el punto objeto y $A'=(0,0,z_i)$ el punto imagen sobre el
eje óptico, ambos con coordenadas axiales positivas $z_0,z_i>0$, y sea
$I=(r,z)$ un punto arbitrario sobre la superficie $\Sigma$, con $r^2=x^2+y^2$
y con distancia al origen dada por
\begin{equation}
\rho^2\;=\;r^2+z^2.
\label{eq:rho_definition_ovoid}
\end{equation}
Los caminos ópticos elementales desde $A$ hasta $I$ y desde $I$ hasta $A'$
son
\begin{align}
\overline{AI}^{\,2}
&= r^2+(z-z_0)^2
= \rho^2-2zz_0+z_0^2,
\label{eq:AI_squared_ovoid}
\\[4pt]
\overline{IA'}^{\,2}
&= r^2+(z-z_i)^2
= \rho^2-2zz_i+z_i^2.
\label{eq:IAp_squared_ovoid}
\end{align}

La condición de Fermat para el estigmatismo exige que el camino óptico total
$L(I)=n_0\overline{AI}+n_i\overline{IA'}$ sea constante sobre $\Sigma$; su
valor se determina evaluando en el vértice $I=(0,0,0)$:
\begin{equation}
L\;=\;n_0\overline{AI}+n_i\overline{IA'}
\;=\;n_0z_0+n_iz_i
\;\equiv\;k.
\label{eq:fermat_principle_for_dioptrique}
\end{equation}

Elevando \eqref{eq:fermat_principle_for_dioptrique} al cuadrado y despejando
el término cruzado,
\begin{equation}
2n_0n_i\,\overline{AI}\,\overline{IA'}
\;=\;
k^2-n_0^2\,\overline{AI}^{\,2}-n_i^2\,\overline{IA'}^{\,2}.
\label{eq:fermat_squared_once}
\end{equation}
Sustituyendo $k^2=(n_0z_0+n_iz_i)^2=n_0^2z_0^2+n_i^2z_i^2+2n_0n_iz_0z_i$ y
las expresiones \eqref{eq:AI_squared_ovoid}--\eqref{eq:IAp_squared_ovoid}, el
lado derecho se reduce a
\begin{equation}
2n_0n_i\,\overline{AI}\,\overline{IA'}
\;=\;
2n_0n_iz_0z_i
-n_0^2(\rho^2-2zz_0)
-n_i^2(\rho^2-2zz_i).
\label{eq:fermat_squared_once_reduced}
\end{equation}

Para eliminar la raíz que aparece en el producto $\overline{AI}\,\overline{IA'}$
se eleva de nuevo al cuadrado \eqref{eq:fermat_squared_once_reduced}. El lado
izquierdo se expande como
\begin{equation}
\begin{aligned}
4n_0^2n_i^2\,\overline{AI}^{\,2}\,\overline{IA'}^{\,2}
&=
4n_0^2n_i^2(\rho^2-2zz_0+z_0^2)(\rho^2-2zz_i+z_i^2)
\\
&=
4n_0^2n_i^2\bigl[(\rho^2-2zz_0)(\rho^2-2zz_i)
\\
&\phantom{=4n_0^2n_i^2\bigl[}
+z_i^2(\rho^2-2zz_0)
+z_0^2(\rho^2-2zz_i)
+z_0^2z_i^2\bigr],
\end{aligned}
\label{eq:LHS_fermat_squared_twice}
\end{equation}
mientras que el lado derecho es
\begin{equation}
\bigl[2n_0n_iz_0z_i-n_0^2(\rho^2-2zz_0)-n_i^2(\rho^2-2zz_i)\bigr]^2.
\label{eq:RHS_fermat_squared_twice}
\end{equation}

Igualando \eqref{eq:LHS_fermat_squared_twice} con
\eqref{eq:RHS_fermat_squared_twice}, cancelando el término común
$4n_0^2n_i^2z_0^2z_i^2$ y agrupando por potencias de las combinaciones
$u\equiv\rho^2-2zz_0$ y $v\equiv\rho^2-2zz_i$, se obtiene la ecuación
polinómica implícita del dioptrio estigmático:
\begin{equation}
\bigl[n_0^2(\rho^2-2zz_0)-n_i^2(\rho^2-2zz_i)\bigr]^2
=
4n_0n_i(n_0z_0+n_iz_i)
\bigl[n_0z_i(\rho^2-2zz_0)+n_iz_0(\rho^2-2zz_i)\bigr].
\label{eq:stigmatic_polynomial}
\end{equation}

La expresión \eqref{eq:stigmatic_polynomial} es una identidad de cuarto grado
en las coordenadas $(\rho,z)$ del punto $I\in\Sigma$. Su forma canónica
—obtenida dividiendo por un factor apropiado en $n_0,n_i,z_0,z_i$ para
la convención axial positiva adoptada— es una cuadrática en $z$ con
coeficientes en $\rho^2$, y se estudia junto con su expansión paraxial en el
apéndice~\ref{apx:ovoids_series}.
